\section{表格型 Q-Learning}

\subsection{Q-Learning 算法}

\begin{importantbox}{Q-Learning 更新规则}
Q-Learning 是一种 off-policy、model-free 的算法，直接学习最优 Q 函数：
$$Q(s, a) \leftarrow Q(s, a) + \alpha \left[r + \gamma \max_{a'} Q(s', a') - Q(s, a)\right]$$
其中：
\begin{itemize}
    \item $\alpha$：学习率
    \item $r + \gamma \max_{a'} Q(s', a')$：TD target（时序差分目标）
    \item $r + \gamma \max_{a'} Q(s', a') - Q(s, a)$：TD error
\end{itemize}
\end{importantbox}

\begin{knowledgebox}{Q-Learning 的 Off-Policy 特性}
Q-Learning 是 off-policy 的，因为它使用 $\max_{a'}$ 来计算 TD target，无论行为策略（用来收集数据的策略）是什么。这意味着我们可以用任意策略（如 $\epsilon$-greedy）来探索，同时学习最优策略。
\end{knowledgebox}

\subsection{探索策略与奖励缩放}

\begin{knowledgebox}{常见探索策略}
\begin{itemize}
    \item $\epsilon$-greedy：大部分时间选最优动作，$1-\epsilon$ 概率探索，适合离散动作空间。
    \item Boltzmann/softmax：通过温度控制随机性，便于把握不同动作的相对价值。
    \item Upper Confidence Bound（UCB）：结合置信区间，在样本不足的状态自动提高探索权重。
\end{itemize}
\end{knowledgebox}

\begin{importantbox}{奖励缩放与非平衡奖励}
Q-Learning 对奖励分布敏感，尺度过大导致 TD error 爆炸，过小又让梯度消失。常用做法包括归一化奖励、基准线减去平均 reward，以及对稀疏奖励做形状工程（dense reward shaping）。在实践中，配套的 instrumentation 应记录各状态-动作的 TD error 和 reward 分布，以便及时调整缩放策略。
\end{importantbox}

\begin{warningbox}{稀疏奖励中的虚假收敛}
当大量状态从未收到正奖励时，agent 可能会在负奖励区域收敛到局部最优策略。建议定期检查 replay buffer 中正负样本的比例，并通过 hindsight experience replay 或多任务辅助奖励增加密度。
\end{warningbox}

\subsection{收敛性条件}

\begin{importantbox}{Q-Learning 的收敛保证}
在表格型设定下（有限状态和动作空间），如果满足以下条件，Q-Learning 保证收敛到 $Q^*$：
\begin{enumerate}
    \item 每个状态--动作对被访问无穷多次
    \item 学习率满足 Robbins-Monro 条件：$\sum_t \alpha_t = \infty$，$\sum_t \alpha_t^2 < \infty$
\end{enumerate}
\end{importantbox}

\subsection{本章小结}

表格型 Q-Learning 通过 TD 更新直接逼近最优 Q 函数，具有理论收敛保证。

